%============================================================
\section{Retirement model with stochastic account realizations}\label{sec:RetirementModelPrices}
%============================================================
The baseline of Section~\ref{sec:RetirementModel} smooths the work--retire margin with directed search and with an endogenous price of the retirement account. That combination takes the place of the Extreme Value Type I taste shocks that usually smooth this margin. This section isolates the account price by shutting search down, and then shows how the two mechanisms work together.

The same idea is described for a durable goods model in Section~\ref{sec:DurableModel}. In the durable model the household searches for a car; contingent on a successful search it observes a stochastic price $p_t$ and then decides whether to buy or to keep. Here the household searches for the option to retire; contingent on receiving that option it observes a stochastic realization of its retirement account and then decides whether to take the draw---retire, absorbing, with that wealth---or to deny it and work another period. Continuing to work is the analog of keeping; claiming the realized account is the analog of adjusting at the realized price. Because $p$ already denotes the exogenous option probability, we write $\phi_t$ for the account realization. Agents who have the option observe $\phi_t$ before choosing and can reject an unattractive draw by remaining at work. The account distribution and claimed wealth are \eqref{eq:ret_price_dist} and \eqref{eq:ret_price_Mtilde}, and the worker's problem is \eqref{eq:ret_full_VW}. Detailed derivations---envelope conditions, Euler equations, the interaction with search, and the EGM algorithm---are in Appendix~\ref{sec:appdx_retirement}.

%------------------------------------------------------------
\subsection{Account-value dispersion on its own}\label{subsec:RetirementPrices}
%------------------------------------------------------------
This subsection shuts search down, in the same spirit as Section~\ref{subsec:ModelBasicPrice}. A worker who receives the option to retire resolves the discrete choice by the hard maximum \eqref{eq:ret_price_discrete}. The comparison itself is contingent on the realized account value; the retiree's problem, once the draw has been accepted, is not. To shut down search we set the exogenous option probability to $p=1$, so the option---and therefore a draw of $\phi_t$---arrives with certainty and optimal effort is identically zero. The account realization is \eqref{eq:ret_price_dist}. Claiming maps book resources into retirement wealth by \eqref{eq:ret_price_Mtilde}, and the value of claiming is \eqref{eq:ret_price_VRet}.

The worker who does not claim solves, with $p=1$,
\begin{align}
	V_t^{W}(M_t)
	&= \max_{0\le c_t\le M_t}
	\left\{
		\log(c_t)-\delta
		+ \beta\,\goth{V}_t^{1}(A_t)
	\right\},
	\label{eq:ret_price_VW}
\end{align}
with $A_t=M_t-c_t$ and $\goth{V}_t^{1}$ the expected value of saving $A_t$ when the option---and a fresh draw of $\phi$---arrives next period. There is no second continuation, because the option arrives with certainty.

A high $\phi_t$ makes claiming more attractive and, if accepted, delivers higher retirement consumption along $c_t^{\Ret}(\cdot)$; a low $\phi_t$ makes claiming less attractive and is denied whenever $V_t^{\Ret}(\phi_t M_t)<V_t^{W}(M_t)$. The model therefore induces state-contingent variation in both the retirement margin and the wealth at which the retiree starts, while leaving the retiree's decision rule unchanged.

The object that matters for the previous period's continuation is the $\phi$-integrated value. With the option arriving for sure it is
\begin{align}
	\bar V_t(M_t)
	\equiv \mathbb{E}_{\phi_t}\big[V_t(M_t,\phi_t)\big]
	= \mathbb{E}_{\phi_t}\Big[
		\max\big\{
			V_t^{\Ret}(\phi_t M_t),\;
			V_t^{W}(M_t)
		\big\}
	\Big],
	\label{eq:ret_price_barV}
\end{align}
which is \eqref{eq:ret_full_barVRW} in the case $p=1$. An Extreme Value Type I taste shock reweights the discrete comparison and leaves both branches' consumption functions unchanged. The account price $\phi_t$ instead changes the argument of the retiree's consumption function. The policy $c_t^{\Ret}(\phi_t M_t)$ is realization-contingent even though the function $c_t^{\Ret}(\cdot)$ itself is not, and the model's reported (ex-ante) consumption function is the $\phi$-weighted average of the realized hard choices. As $\sigma_\phi\to 0$, $\phi_t\to 1$ almost surely and, with $p=1$, the discrete choice is a hard maximum at book resources.

Two differences from the durable-price extension are worth flagging. First, retirement is self-absorbing, so a household that accepts $\phi_t$ cannot return to work next period to wait for a better draw; the analog of keeping is available only before claiming. Second, because $\phi_t$ rescales the state rather than the first-order condition, the retiree is solved once, and each $\phi$-node is only an evaluation point---cheaper than re-solving the durable adjuster at every $p_k$. The smoothing mechanism is the same: integrating \eqref{eq:ret_price_discrete} over $\phi_t$ replaces a single wealth threshold $M^\ast$ at which $V^{\Ret}(M^\ast)=V^{W}(M^\ast)$ with a family of thresholds $M^\ast(\phi)$, so the ex-ante claiming probability is interior and the expected marginal value that enters last period's Euler equation is smoothed. We derive the envelope conditions, the terminal-period solution, and the EGM algorithm for this case in Appendix~\ref{subsec:appdx_ret_price}.

%------------------------------------------------------------
\subsection{Search and the account price together}\label{subsec:RetirementSearchPrices}
%------------------------------------------------------------
Putting search back in recovers the baseline of Section~\ref{sec:RetirementModel}. The household chooses search effort; contingent on a successful search it receives the option to retire and a draw of the retirement account; it then takes the draw or works another period. The exogenous option probability $p\in[0,1)$ is restored, so a worker who exerts no effort receives a draw with probability $p$ and a worker who exerts maximal effort receives one with certainty. The worker's problem, the two gothic values, the effort rule, and the $\phi$-integrated value are \eqref{eq:ret_full_VW}, \eqref{eq:ret_full_gothV1}, \eqref{eq:ret_full_gothV2}, \eqref{eq:ret_full_effort}, and \eqref{eq:ret_full_barVRW}.

The two frictions play complementary roles, as in Section~\ref{subsec:ModelBasicFull}. Search governs whether a draw arrives. The account price $\phi_t$ is the content of the draw, and the hard maximum \eqref{eq:ret_full_discrete} is the decision to take it or to work another period. Search alone cannot cherry-pick a good account realization, because effort is chosen before $\phi_{t+1}$ is observed. Account dispersion on its own, the case $p=1$ above, cannot smooth the problem of a worker who did not receive the option, because that agent cannot claim at any $\phi$. Together they smooth both the arrival of the option and the terms on which it is exercised. With $\sigma_\phi\to 0$ the baseline collapses to search with a hard maximum at book resources. With $p=1$ it collapses to \eqref{eq:ret_price_VW}. We derive the combined envelope conditions, the worker's Euler equation, the terminal-period solution, and the EGM algorithm in Appendix~\ref{subsec:appdx_ret_full}.

\subsection{Results}
Introducing price dispersion smooths the problem numerical with two features which figure~\ref{fig:fig8} illustrates. The left panel shows the consumption function for larger price dispersion. A higher price dispersion means that the agent could receive a much larger retirement wealth amount or a smaller amount. Therefore, agents who receive a low draw might want to wait another period and hope for a better draw in the future. As a result, larger price dispersion leads to a smoother consumption function. Numerically, it matters how many price points we are considering. The right panel shows that with only 5 points (red), we increase the number of discontinuities, but it is not enough to smooth the full consumption function. With a larger number of points and th assumption of a large number of agents, the combined consumption function consists of the average of many functions in which some agents accept and others do not. Note that for a large number of shock points -- eg 151 -- the smoothing does not change much compared to an even larger number --eg 351.


%---- Figure 8 ----
\begin{figure}[htbp]
	\centering
	\begin{subfigure}{0.48\textwidth}
		\includegraphics[width=\linewidth]{VaryPhi.pdf}
		\caption{Varying the price dispersion of retirement wealth.}
	\end{subfigure}
	\hfill
	\begin{subfigure}{0.48\textwidth}
		\includegraphics[width=\linewidth]{VaryPhiCount.pdf}
		\caption{Varying the number of price dispersion shock points.}
	\end{subfigure}
	\caption{The plots show the impact of increasing the smoothness through larger standard deviation of price dispersion
	(left panel) and increasing the number of price shock points (right panel).}
	\label{fig:fig8}
\end{figure}

Additionally, we can combine search with price dispersion. Since agents only receive a draw if they are searching, combining both methods amplifies the smoothing. Figure~\ref{fig:fig9} shows on the left that increasing search while keeping price dispersion small leads to a powerful smoothing mechanism. The right panel shows the opposite. When we keep search intensity as given, but increase price dispersion, the consumption function becomes very smooth and differentiable.

%---- Figure 9 ----
\begin{figure}[htbp]
	\centering
	\begin{subfigure}{0.48\textwidth}
		\includegraphics[width=\linewidth]{SearchPhi_VaryS.pdf}
		\caption{Varying the search intensity given price dispersion of retirement wealth.}
	\end{subfigure}
	\hfill
	\begin{subfigure}{0.48\textwidth}
		\includegraphics[width=\linewidth]{SearchPhi_VaryPhi.pdf}
		\caption{Varying price dispersion of retirement wealth given search intensity.}
	\end{subfigure}
	\caption{The plots show the impact of smoothness through the combination of search and price dispersion. Left panel keeps price dispersion fixed and varies search intensity and right panel keeps search intensity fixed and varies price dispersion.}
	\label{fig:fig9}
\end{figure}